% Directly included technical appendix to compact_arithmetic20.tex.
% The complete local geometric derivations are an explicit internal supplement.
\appendix
\section{A rank 20 improvement in eight ground fields}\label{app:eight}

Put $K_m=\Q(2\cos(2\pi/m))$ and let
\[
 \mathcal K=\{\Q(\sqrt3),\Q(\sqrt6),K_9,K_{11},K_{15},K_{16},K_{20},K_{24}\},
 \qquad K_{24}=\Q(\sqrt2,\sqrt3).
\]
Here and below a statement about $K$ means equality with the actual
ground field, rather than containment in a displayed overfield.
\begin{theorem}\label{thm:eight19}
A compact arithmetic hyperbolic Coxeter polytope whose ground field
belongs to $\mathcal K$ has dimension at most 19.
\end{theorem}
Theorem~\ref{thm:main} already excludes dimensions at least 21. This
appendix excludes dimension 20, without extrapolating its rank 20
calculations to other dimensions. It does not improve the all-field
ceiling 20. In particular it makes no ceiling 19 assertion for
$\Q(\sqrt2),\Q(\sqrt5),\Q(\sqrt2,\sqrt5)$ or $K_7$.

We give the complete composition of a new half-refund rule with the
previous charge, its triangle and quadrilateral obligations, and the
strict counting contradiction. The six local quadrilateral exclusions
are imported \emph{internal geometric theorems}, with precise statements
below. Their complete proofs form the named reproducibility supplement;
they are not published inputs or consequences of the numerical rank 21
certificate. Thus this appendix is a complete composition proof, not a
self-contained replacement for every local geometric derivation.
The machine-readable ledger is
\path{manuscripts/supplements/compact_arithmetic_eight_field19_sources_final_20261019.json}.
It records exact source bytes and roles, including the required
spectator-$E_8$ coverage correction.

\subsection{Inputs, normalization and the previous charge}
All local roots in this appendix are unit outward normals. Distinct
normals have nonpositive scalar product. A positive-definite root set
defines an actual face, and a positive affine set is impossible, by
Lemma~\ref{lem:projection}. A compact acute chamber is simple. We use
the finite and affine Coxeter classification, arithmetic cyclic-product
integrality and strict positivity at nonphysical embeddings, and the
nonnegative symmetric $h$-vector theorem as external structural inputs.
Arithmeticity of reflection groups in the relevant dimensions is of
the simplest type, as in the main field reduction.

One arithmetic observation will recur. Add an actual wall $c$ to the
\emph{entire} positive rank 20 vertex diagram. At the old simple vertex
the new wall inequality is strict, so the resulting 21 normals are
independent and their physical Gram matrix is Lorentzian. The projection
cost $q_c$ against the old vertex is therefore greater than 1 physically
and less than 1 at every nonphysical embedding. If $q_c$ is a rational
expression in a cyclic scalar $z\in K$, then
\begin{equation}\label{eq:eight-field}
 K=\Q(z).
\end{equation}
Indeed, if this were a proper subfield, another embedding of the totally
real field $K$ extending its physical embedding would repeat $q_c>1$.
If the cost is rational, the same argument forces $K=\Q$.
This observation concerns the whole old vertex; it is not valid for
an arbitrary embedded component with unexamined outside contacts.

The initial charge on a marked pair is one when its distance in its
finite component is at most two, and zero otherwise. Add one more
when the pair is an \emph{entire} component $I_2(m)$ with $m\ge6$.
Call this the original charge. Its triangle payments are at least 2
and its quadrilateral payments at least 1 over $\mathcal K$, by the
internal high-label theorem. Its complete finite-domain proof is an
input, not a new assertion about an arbitrary arithmetic field.

We retain the previously proved whole-$A$, whole-$D$, and $B$ reductions
and their positive coupons. The inherited internal charge theorem used
here has the following exact scope and consequences.
\begin{proposition}[Previous charge interface]\label{prop:eight-old}
There is a nonnegative charge on every actual rank 20 vertex over
$\mathcal K$, obtained from the original charge by the reviewed
whole-$A/D/B$ rules, with every triangle paying at least 2 and every
quadrilateral at least 1. Its only vertex costs above 35 are
\[
 D_r+D_{20-r}\quad(10\le r\le16),\qquad
 D_{12}+E_8,\quad D_{13}+E_7,\quad D_{14}+E_6,
\]
each with cost 36. It leaves both fork/short-leaf units in these forests
unchanged. Its other exact bounds are
\[
\begin{array}{c|c}
 \text{entire vertex forest}&\text{cost bound}\\\hline
 A_{20},B_{20},D_{20}&35\\
 D_{19}+A_1,\ D_{18}+I_2(m),\ m\ge6&35\\
 B_{19}+A_1&139/4\\
 \text{at least three components}&34.
\end{array}
\]
The only positive coupon with at least three components is the
$B_{18}+2A_1$ quarter, whose complete cost is $133/4$.
\end{proposition}
This is the composed prior theorem, not an application of a whole-$D$
refund to a $D$ factor with spectators. Its inputs are the whole-$A$
two-end refund, whole-$D$ three-unit refund, their mixed-consumer
theorem, the $B$ simple-end and balanced-quarter rules, the $B_{20}$
far-four quarter, and the spectator rules for $B_{19}+A_1$,
$D_{19}+A_1$ and $D_{18}+I_2(m\ge6)$. All are identified in the ledger.

For clarity, the finite-forest part needs no new enumeration. A
connected rank-$r$ component has original cost $2r-2$ for $D/E$,
for an entire high-label $I_2$, and for $A_1$ (cost zero); every other
nontrivial component has cost $2r-3$. Indeed a tree contributes its
edge count plus $\sum_v\binom{\deg v}{2}$ distance-two pairs. Thus
a forest with $k$ components has original cost at most $40-2k$.
For $k\ge3$ the only positive coupon is the stated $133/4<34$ case.
For $k=2$ any nontrivial path component without a high-$I_2$ addition
gives cost at most 35, apart from the separately bounded $B_{19}+A_1$.
Two exceptional $E$ components cannot fill rank 20. The only remaining
cost 36 cases are $D+D$, $D+E$, $D_{19}+A_1$ and
$D_{18}+I_2(m\ge6)$; the last two have their own prior half-refunds.
For $k=1$ only $A_{20},B_{20},D_{20}$ occur, already bounded by 35.
This proves the stated frontier from the earlier consumer theorems.

\subsection{The new rule and its complete support}
At one of the ten cost 36 forests select the displayed larger $D_r$
factor. At $D_{10}+D_{10}$ select either factor independently at each
vertex. Write $a$ for its fork, $b,e$ for its short leaves and
\[
 a-l_1-\cdots-l_L,\qquad N=r-2\ge8,\quad L=N-1,
\]
for its long arm. Let $T_s$, $r+s=20$, be the entire other $D/E$
component. Reduce the two weights on $(a,b)$ and $(a,e)$ from 1 to
$1/2$. Leave the short-leaf pair $(b,e)$ and every other angle unchanged.

The complete white support of a face carrying the pair $(a,b)$ is
\begin{equation}\label{eq:eight-white}
 W=A_1(e)+A_L+T_s.
\end{equation}
Its branched component $T_s$ distinguishes it from all previous
decreasing supports: the whole-$A/D$ roles have $A_{18}$ or
$A_{17}+A_1$, the $B$ roles have $A_{18}$ or a white bond of label 4,
$D_{19}+A_1$ has $A_{16}+2A_1$, and the high-$I_2$ role has
$A_{15}+A_1+I_2(m\ge6)$. The positive $B$ coupons have white support
$B_{18}$. None equals~\eqref{eq:eight-white}. Consequently no previous
decrease occurs at any other corner of this face, and an original
positive corner still contributes at least one unless it is a new
half-source.

Among new roles the unique branched white component fixes $T_s$, and
the path rank fixes $r$. Thus simultaneous new reductions on this face
belong to the same family. In the tie case the whole white $D_{10}$
already occupies its entire factor, so its roots cannot secretly supply
a black selected pair. Omitting a reduction because a different factor
was selected at another vertex can only increase payment. No consistent
global choice of the selected factor is needed.

\subsection{Every triangle retains two units}\label{app:eight-triangles}
Let $a,b,c$ be the sides of an actual triangle with a new source at
$(a,b)$. Both other corners are finite. In particular
\[
 W+a=A_{r-1}+T_s,\qquad A_{r-1}=e-a-l_1-\cdots-l_L.
\]
We first examine contacts to the whole spectator. A finite diagram is
a forest, so $c$ touches at most one vertex in each connected component.
If $c$ touches only $T_s$, a nonzero $b.c$ and the spectator contact
must both be simple: a finite extension retaining a $D/E$ fork has only
simple bonds. The whole old-plus-$c$ Gram data are then rational,
contrary to~\eqref{eq:eight-field}. If $b.c=0$ that whole rank 21
diagram is positive definite, impossible at positive index 20.
If $c$ also touches the long path, the whole finite corner of rank 20
is connected and contains $T_s$'s fork, so it must be $D_{20}$.
The spectator is $D_s$ and $c$ joins its long endpoint to a path
endpoint simply. A nonzero $b.c$ would make a second fork in the
other finite corner. At $b.c=0$ the whole old-plus-$c$ cost is again
rational and physically greater than 1, forcing $K=\Q$. Hence
\begin{equation}\label{eq:eight-triangle-seal}
 c\perp T_s
\end{equation}
in every possible triangle. Orthogonality has been proved after
examining spectator contacts, not assumed at the outset.

Finite extension of the whole path $A_{r-1}$, of rank at least 9,
now leaves five cases: no path contact; a simple or label 4 contact at
either endpoint $e,l_L$; or a simple contact at either next-to-end
vertex $a,l_{L-1}$. Two path contacts create a cycle, and the
exceptional finite diagrams cannot extend this path at its rank.
Put $b.c=-\beta$. We show that each triangle remains paid even
if its new fork units were deleted completely.

\paragraph{Near endpoint $e$.}
If $c.e$ has label 3 or4 and $\beta>0$, the finite path $e-c-b$
allows ordered labels $(3,3),(3,4),(3,5),(4,3)$. The marks $(a,c)$
and $(b,c)$ have respectively distance two and one. Neither is a new
source, so they retain one unit each. At $\beta=0$, label 3 gives
the affine $E_8$ tree with arms $1,2,5$ at $a$, using
$a,b,e,c,l_1,\ldots,l_5$. Label 4 gives the affine block on
$a,b,l_1,e,c$ with positive null coefficients $2$ at $a,e$,
$1$ at $b,l_1$, and $\sqrt2$ at $c$. Both are forbidden.

\paragraph{Far endpoint $l_L$.}
At $\beta=0$, a simple contact makes the whole rank 21 diagram
$D_{r+1}+T_s$ positive definite. A label 4 contact gives a positive null
vector with coefficients $1$ on $a$ and the long arm, $1/2$ on
$b,e$, and $1/\sqrt2$ on $c$. At $\beta>0$, finite classification
of the opposite path forces the far contact simple and
$\beta=1/2$ or $\sqrt2/2$. The first makes an affine simple cycle.
For the second, the full old-plus-$c$ projection cost is
\[
 q_c=r\beta^2/2+\beta+1/2.
\]
The actual cycle puts $\sqrt2$ in $K$. Extend its sign-changing
embedding to $K$; it gives
$q_c'=r/4-\sqrt2/2+1/2>1$, contradicting conjugate positivity.

\paragraph{Far next-to-end $l_{L-1}$.}
The contact is simple. At $\beta=0$ there is an affine $D$
vector: coefficients $1$ on $a,l_1,\ldots,l_{L-1}$ and $1/2$
on $b,e,l_L,c$. At $\beta>0$ the other finite corner would
have fork arms $1,2,L-2$ and rank $r-1\ge9$, which is not a
finite diagram.

\paragraph{Near next-to-end $a$.}
Here $a.c=-1/2$ and $c\perp W$. At $\beta=0$ the star with
center $a$ and leaves $b,e,l_1,c$ is affine $D_4$. Otherwise
the two corners with pairs $(a,b),(a,c)$ are new fork roles, and
the third is $W+I_2(m)$, where $\beta=\cos(\pi/m)$.
The actual black cycle makes $2\beta$ integral in $K$, and
\[
 q_c=\frac{r\beta^2+2(r-2)\beta+r-2}{2},\qquad
 q_c-1=\frac{(\beta+1)(r\beta+r-4)}2.
\]
Strict positivity gives, at each nonphysical embedding,
\begin{equation}\label{eq:eight-near-window}
 -1<\beta_i<-1+4/r\le-3/5,
\end{equation}
and the whole-star argument gives $K=\Q(\beta)$.
Every field in $\mathcal K$ has degree at most 5. The elementary
totient condition $\varphi(2m)\le10$ gives exactly
\[
 m=3,4,5,6,7,8,9,10,11,12,15.
\]
The prime divisors of $2m$ must be among $2,3,5,7,11$ and
multiplying their Euler factors gives this complete list.
Label 3 is affine $A_2$; labels 4,5 give the excluded exact fields
$\Q(\sqrt2),\Q(\sqrt5)$. For $m=7,8,10,11$ use conjugate
$\cos(3\pi/m)$, for $m=9,12$ use $\cos(5\pi/m)$, and for
$m=15$ use $\cos(7\pi/m)$. These are genuine nonphysical
conjugates and violate~\eqref{eq:eight-near-window}: all are
positive except $\cos(5\pi/9)>-1/2$. Only $m=6$ remains.
Its third corner is an \emph{entire} $I_2(6)$ and retains two
units, paying the triangle even if both fork units were removed.

\paragraph{No path contact.}
Now $c\perp a,W$. At $\beta=0$ the whole rank 21 set is positive
definite. Otherwise the integral two-cycle
$z=4\beta^2=2+2\cos(2\pi/m)$ satisfies
\[
 q_c=rz/8,\qquad K=\Q(z),\qquad 0<z_i<8/r\le4/5.
\]
The complete condition $\varphi(m)\le10$ gives
\[
 m=3,4,5,6,7,8,9,10,11,12,14,15,16,18,20,22,24,30.
\]
Labels 3,4,6 give $\Q$; labels 5,8 give
$\Q(\sqrt5),\Q(\sqrt2)$. For all remaining labels except 12,
the following nonphysical residues are coprime to $m$, not $\pm1$,
and have $k/m<1/3$:
\[
\begin{array}{c|rrrrrrrrrrrr}
m&7&9&10&11&14&15&16&18&20&22&24&30\\
k&2&2&3&2&3&2&3&5&3&3&5&7.
\end{array}
\]
Their $z_i=2+2\cos(2k\pi/m)>1$ violate the window. This keeps
label 20 explicitly; it does not assume $\beta$ belongs to its
square field. Only $m=12$ remains, with exact field
$\Q(\sqrt3)$. The entire $I_2(12)$ third corner retains two
units. Thus every triangle, including all simultaneous new reductions,
pays at least 2.

\subsection{Quadrilaterals: complete reduction to six local lemmas}
Two new half-sources already pay one. If a quadrilateral has exactly
one new half-source and another original positive corner, that other
corner retains at least one by the white-support separation above.
We must therefore exclude just a quadrilateral $a,b,c,d$ with its
source at $(a,b)$ and three other \emph{original} zero corners.
The zeros imply
\begin{equation}\label{eq:eight-zero}
 a.d=b.c=c.d=0,\qquad d.e=d.l_1=0.
\end{equation}
The latter equalities hold because a contact to either root, both
simply adjacent to $a$, would give a distance-two unit at $(a,d)$.
Put $b.d=-\beta$ and $a.c=-x$. Actual acute Coxeter walls give
$x,\beta=0$, a native finite cosine, or a value at least 1.
We do not infer dotted ambient products merely from polygon opposition.

Finite extension of the \emph{entire} $A_{r-1}+T_s$ at $W+a+d$
and~\eqref{eq:eight-zero} give precisely the first-$d$ inventory
in Table~\ref{tab:eight-d}. A near attachment is excluded by the
zero at $(a,d)$. Two contacts in one old component create a cycle.
A path-plus-spectator connected finite diagram has rank 20 and a
fork, so it must be $D_{20}$: the bridge is simple and uses a
$D$ long endpoint. Neither an $E$ spectator nor a second fork nor
a bond of label 4 in that connected component is finite. The old path has
rank at least 9, so no exceptional internal attachment is omitted.
\begin{table}[ht]
\centering
\begin{tabular}{cl}\toprule
Row&Contacts of $d$ to the complete white graph\\\midrule
I&None\\
F&Simple far endpoint $l_L$\\
B&Far endpoint $l_L$ with label 4\\
T&Simple far next-to-end $l_{L-1}$\\
J&One simple contact to $T_s$, no path contact\\
Bridge&Simple $l_L$ and a long endpoint of $T_s=D_s$\\\bottomrule
\end{tabular}
\caption{The complete first-wall inventory for three original zero corners.}
\label{tab:eight-d}
\end{table}

Here are the exact local geometric inputs. Each statement has the
entire old forest and all hypotheses~\eqref{eq:eight-zero}; it
excludes every actual finite continuation $W+c+d$, not just an
embedded classical subdiagram.
\begin{proposition}[Internal six-row exclusion]\label{prop:eight-six}
For $10\le r\le16$, $r+s=20$, $T_s=D_s,E_6,E_7,E_8$ and
$K\in\mathcal K$, none of the six rows in
Table~\ref{tab:eight-d} can be an actual compact quadrilateral
with three original zero corners. Both finite and nonfinite ambient
opposite products are included. The J statement includes the two
$A_7$-to-$E_8$ continuations at $N=8$.
\end{proposition}
This proposition is an imported internal theorem assembled from the
local sources below. We give its exact proof interface and principal
calculations to make its coverage assessable. The source ledger provides
the full individual derivations; accepting a smaller-rank extension of
an old $D_{19}/D_{20}$ estimate is not one of the premises.

\paragraph{The common determinant identity.}
Project off the positive white block and write
\[
 \begin{pmatrix}
 s_0&-1/2&-X&-\rho\\
 -1/2&1&0&-\beta\\
 -X&0&t&-\nu\\
 -\rho&-\beta&-\nu&u
 \end{pmatrix},\qquad
 s_0=\frac{N+2}{4N},\quad t,u>0,\quad tu>\nu^2.
\]
There are 22 actual roots in Lorentz dimension 21, with a whole positive
corner of rank 20, so this matrix has determinant zero. When
$\Delta=\beta^2-u>0$, eliminating the $b,d$ block gives
\begin{equation}\label{eq:eight-schur}
 X=\frac{\nu(\beta/2+\rho)}\Delta+
 \sqrt{\left(t+\frac{\nu^2}\Delta\right)
 \left(s_0+\frac{\rho\beta+\rho^2+u/4}\Delta\right)}.
\end{equation}
The other algebraic root is negative: the constant term of the
quadratic has the sign forced by $s_0>1/4$. Thus the positive root
is determined geometrically. For $\beta\ge1$ its right side is
nonincreasing: the relevant derivative numerators are
\[
 -\beta^2/2-2\rho\beta-u/2,\qquad
 -\rho\beta^2-\rho u-2(\rho^2+u/4)\beta,
\]
both negative. Replacing $\beta$ by 1 is therefore an upper estimate.

\paragraph{Row I: actual dottedness and the complete contact split.}
Here $b,d\perp W$ and are also perpendicular to the projected $c$
side. The actual compact polygon has a positive-length $c$ edge;
its two adjacent side lines $b,d$ are perpendicular to that same
line at distinct points. Their outward product is
$-\cosh\ell$, so $\beta>1$ is derived from actual geometry.
Put $\delta=4\beta^2-4\in\mathcal O_K$. The identity becomes
\begin{equation}\label{eq:eight-I}
 \left(\frac{(x+\kappa)^2}{t}-s_0\right)\delta=1,
 \qquad q_c=1-t+2N(x+\kappa)^2.
\end{equation}
The complete $c$ split is: no long-path or spectator contact, with
arbitrary finite $e$ contact; a long-path contact with every allowable
$e$ contact; spectator only; spectator plus a simple $e$ contact; or a path/spectator
bridge. The no-spectator branches include both $N=8$ exceptional
$E_8$ attachments. Their full no-long, finite-long and dotted-long
unit arguments are retained as separate internal lemmas in the
supplement. No bounded algebraic-integer search is an input.

For a spectator contact the residual squares after its projection are
\[
\begin{array}{c|c|c}
T_s&\text{attachment}&t_T\\\hline
D_s&\text{long endpoint (any leaf for }D_4)&1/2\\
D_5,D_6,D_7&\text{short leaf}&3/8,1/4,1/8\\
E_6&\text{length-two arm endpoint}&1/3\\
E_7&\text{length-three arm endpoint}&1/4\\
E_8&\text{none}&\text{--}.
\end{array}
\]
This follows from finite $D/E$ arm lengths and unit-Gram determinant
ratios. A spectator-plus-$e$ contact must be simple and leaves only
$(t,N)=(1/4,8\ldots14),(1/8,13),(1/12,12)$.
With the integral actual-cycle scalar $A=2x$ and $V=2A+1$,
equation~\eqref{eq:eight-I} respectively produces
\[
\begin{array}{c|c|c}
t& D_*&D_*\delta\\\hline
1/4&2NA(A+1)-1&2N\\
1/8&8NA(A+1)+N-2\quad(N=13)&52\\
1/12&72A(A+1)+11\quad(N=12)&24.
\end{array}
\]
In every row $D_*$ is coprime to its displayed divisor, hence is a
unit. Either $x=0$ contradicts $\delta>0$ directly, or its physical
absolute value exceeds 1; the full-star inequality gives absolute
value greater than 1 at every other conjugate too. This contradicts
its unit norm. For a path/spectator bridge the whole finite component
is $D_{18}$, its residual is $1/(2N)$, and its two endpoint roles
give an odd unit $D_*$ satisfying $D_*\delta=4$; the same conjugate
norm contradiction applies. The spectator-only row is settled by
the coupled integral-square norm argument. These are the five
complete contact alternatives in the named I composition lemma.

\paragraph{Rows F and T at finite $\beta$.}
The integral cycle $2\beta$ belongs to $K$. The full old-star costs are
\[
 q_F=r\beta^2/2+\beta+1/2,\qquad
 q_T=r\beta^2/2+2\beta+1.
\]
They give $K=\Q(\beta)$ and the common nonphysical window
$-1/2<\beta_i<1/4$. The complete finite list is the
$\varphi(2m)\le10$ list above. Labels 3,5 force excluded fields.
For the other labels a forbidden nonphysical residue is
\[
\begin{array}{c|rrrrrrrrr}
m&4&6&7&8&9&10&11&12&15\\
k&3&5&5&3&7&3&3&5&11.
\end{array}
\]
Negative entries have angle greater than $2\pi/3$; the positive
exceptions exceed $1/4$, including
\[
 \cos(5\pi/12)=(\sqrt6-\sqrt2)/4>1/4.
\]
Also $\beta=0$ fails the physical cost $q_d>1$.

\paragraph{Row B at finite $\beta$: the weighted-square exception.}
The actual integral scalar is $A=2\sqrt2\beta$, not necessarily
$\beta$. Its full cost is $rA^2/16+A/2+1$, giving
\[
 K=\Q(A),\qquad -8/r<A_i<0.
\]
If $\Q(A^2)$ were proper, the sign automorphism $A\mapsto-A$
would give the forbidden conjugate $-A_0\le-\sqrt2$.
Thus $K=\Q(A^2)=\Q(\cos(2\pi/m))$ and the complete
$\varphi(m)\le10$ list applies, including $m=20$.
All labels except 12 fail the necessary nonphysical square bound
$\sigma(A^2)<64/r^2\le16/25$: for $m=5,8$ the values
$3-\sqrt5>3/4$ and $4-2\sqrt2>1$ suffice; for the other
nonrational labels use the residues $k/m<1/3$ listed above,
giving $4+4\cos(2k\pi/m)>2$.
For $m=12$, $A=1+\sqrt3$, so its conjugate window leaves only
\[
 r=10,\quad T_s=D_{10},\quad \beta=\cos(\pi/12),\quad K=\Q(\sqrt3).
\]
This necessary profile is not discarded. In it the complete opposite
corner is $A_1+B_8+D_{10}$: $c$ has no path contact, or a simple
near-end contact $l_1$ with $e$ absent or simple. Eliminating $b,d$
gives $R_0=(2+\sqrt3)/8<1/2$. For no path contact,
\[
 (x+\gamma/2)^2=R_0(t_T-\gamma^2),\qquad \gamma=-c.e.
\]
The native gap leaves $x=0,1/2$. A spectator contact makes $t_T$
rational, whereas the remaining possible equalities require
$4-2\sqrt3$ or $1-\sqrt3/2$ as residuals. Without a spectator,
the last $x=0$ case requires
$4\gamma^2=(20+8\sqrt3)/13$, whose trace $40/13$ is not integral.
For the near-end contact the exact identity is
\[
 (x+3/8+\gamma/2)^2
 =R_0\left((12+\sqrt3)/24-\gamma^2\right),
 \qquad \gamma=0,1/2.
\]
A positive $x$ makes the left side at least $49/64$, too large;
at $x=0$ the two right numerators are respectively
$27+14\sqrt3$ and $15+8\sqrt3$ over 192, different from the
left numerators 27 and 75. This closes the genuine finite exception.

\paragraph{Rows F, B and T for $\beta\ge1$.}
For F the opposite path is $A_N$, for B it is $B_N$, and for T
it is $D_N$. Their complete continuations are constrained by the
whole finite corner, not by an assumed spectator seal. Equation
\eqref{eq:eight-schur} has the following exact data and bounds:
\[
\begin{array}{c|c|c|c}
\text{row}&u&\rho&Q(1)\text{ bound}\\\hline
F&(N+1)/(2N)&1/(2N)&(N+3)/(2(N-1))\le11/14\\
B&1/N&\sqrt2/(2N)&(5+\sqrt2)/14<1/2\\
T&2/N&1/N&(N+6)/(4(N-2))\le7/12,
\end{array}
\]
where $Q(\beta)$ is the second square-root factor. If $c$ has no
path contact, $(x+\gamma/2)^2/t\le Q(1)$. The spectator table
gives $t\le1/2$, or $t\le1/4$ with a simple $e$ contact. Native gaps,
the induced affine $D/E/B$ witnesses, and the full-star exact fields
$\Q$ or $\Q(\sqrt5)$ exclude these cases.

For B and T the only path contact is simple at $l_1$, with
$\gamma=0,1/2$ and no spectator contact. In T it gives
$X\le11/16$, while a positive $x$ gives $X\ge15/16$.
In B the corresponding upper bound is
$X<1/8+\sqrt{2/7}<7/8<15/16$. Thus $x=0$, giving affine
$D_5$ at $\gamma=0$ or the simple affine four-cycle at $\gamma=1/2$.

For F the full path-contact list is a simple or label 4 near endpoint,
simple near-next or far-next, or a simple near-end/spectator bridge.
The respective estimates force $x=0$ with an affine witness, or
contradict even its zero-contact lower bound. For example the
simple near-end case has
\[
 X\le9/112+\sqrt{22}/7<15/16,
 \qquad X\ge15/16\quad\text{if }x>0;
\]
the near-four and bridge upper bounds are respectively
$9\sqrt2/112+\sqrt{11}/(7\sqrt2)<7\sqrt2/16$ and
$9/112+\sqrt{11}/14<7/16$, below their lower bounds at $x=0$.
The sole shared-far-end alternative instead gives
\[
 x\le1/2+\sqrt{(N+3)/(2(N-1))}<3/2.
\]
Its integral scalar $A=2x$ generates $K$ and has all other
conjugates in $(-1/2,1/4)$, since
$q_c=NA^2/2+A+1/2$. Finite native $x$ fails the complete
cosine list. For $x\ge1$ and degree at least 3,
$0<|\operatorname{Norm}A|<3(1/2)^{d-1}<1$.
In degree two a positive other conjugate also has norm below 1;
a negative one forces norm $-1$ and trace 2, hence
$A=1+\sqrt2$ and the excluded exact field $\Q(\sqrt2)$.
Degree one forces $\Q$. This is the all-height far-end argument,
not a height-box census.

\paragraph{Row J, with the small exceptional continuations retained.}
A simple spectator contact has rational residual $u\le1/2$.
The full-star integral square $z=4\beta^2$ generates $K$ and
has $0<z_i<8u/r\le2/5$. Integer norms and traces exclude every
finite native $\beta$ except
\[
 \beta=\cos(\pi/12),\quad K=\Q(\sqrt3),\quad
 T_s=D_s\text{ at its long endpoint},\quad 10\le r\le14.
\]
This exception is kept in the local theorem. If both $c,d$ touch
the spectator, the complete joint finite table is
$D_4\to E_6,D_5\to E_7,D_6\to E_8$, with projected cross
$\nu=1/4$ and no path contact. Their remaining constants satisfy
$s_0\le7/24$, $u,t_T\le1/2$; an $e$ contact is absent or simple.
For $\beta\ge1$, equation~\eqref{eq:eight-schur} gives
$X<5/6$ when $\gamma=0$, and $x<\sqrt{13}/8<1/2$ when
$\gamma=1/2$. The remaining native labels have rational or
$\Q(\sqrt5)$ full-star costs and are excluded. At the finite
exception only the $D_6\to E_8$ joint row remains; its exact
comparison gives $X<4/5$, leaving $x=0,1/2,\sqrt2/2$, each
excluded by its full-star cost.

If $c\perp T_s$, then $\nu=0$ and
\[
 X^2=tR,\qquad R=s_0+\frac{u}{4(\beta^2-u)}.
\]
For $\beta\ge1$, $s_0<R\le(N+1)/(2N)\le9/16$; in the
finite exception $R=s_0+\sqrt3/6$, strictly between
$(N+1)/(2N)$ and $5/8$. The complete classical path/end/$e$
inventory is excluded by the native gaps, affine witnesses and
rational full-star costs in the J source. There are also two
\emph{additional} continuations when $N=8$: $c$ can attach
simply to $l_3$ or $l_5$, making $E_8$ from the white $A_7$.
No $e$ contact is then possible. They have
\[
 t=1/16,\quad X=x+\kappa,\qquad
 \kappa=5/16\ (l_3),\quad 3/16\ (l_5).
\]
For $\beta\ge1$, $X\le3/16$ excludes $l_3$; $l_5$ forces
$x=0$ and the induced affine $D$ tree with its two forks at
$a,l_5$. In the finite exception $X<1/5$, so again $l_3$
fails, while $l_5$ forces $x=0$ and the impossible rational
value $R=9/16$. These are the mandatory $E_8$ addendum cases.

\paragraph{The bridge row.}
The whole opposite white factor is $D_{18}+A_1(e)$. Thus $c$
either misses $D_{18}$ and has any finite $e$ bond, or lengthens
its long endpoint $l_1$ simply and has $e$ absent or simple.
The integral $V=2\beta$ has $K=\Q(V)$ and
$-4/r<V_i<0$. Finite $\beta$ gives
$0<|\operatorname{Norm}V|<2(4/r)^{d-1}<1$ for $d\ge2$;
degree one forces $\Q$. For $\beta\ge1$, the determinant
data are $u=\rho=1/(2N)$ and
\[
 Q(1)=\frac{N+4}{2(2N-1)}\le2/5.
\]
No $D_{18}$ contact leaves the native labels 3,4, excluded by
the usual affine $D/E/B$ witnesses. A simple $l_1$ contact gives
$X\le3/80+\sqrt{17/75}<43/80$, below the $15/16$ lower
bound when $x>0$. At $x=0$ its absent/simple $e$ bond gives
affine $D_5$ or the simple affine four-cycle.

Every affine witness above is principal on actual outward roots.
Its positive null sum is nonzero by strict interior pairing and
is feasible against \emph{all} other walls by acuteness, including
any outside spectator walls. Additional contacts cannot remove
the ideal ray. The complete local proofs supply all finite-corner
inventories and the subsidiary native comparisons summarized here.
With Proposition~\ref{prop:eight-six}, the one-half-source/three-zero
quadrilateral is impossible. Every affected quadrilateral therefore
retains payment at least 1.

\subsection{Vertex bound, strict contradiction and verification}
The new charge remains nonnegative. Its triangles pay 2 by
Section~\ref{app:eight-triangles}; its quadrilaterals pay 1 by the
preceding complete consumer argument. Every target forest loses
exactly one: its two distinct original units become two halves.
No positive coupon occurs on a target, and no unproved whole-factor
deletion is applied through spectators. Every other forest is unchanged.
Proposition~\ref{prop:eight-old} consequently gives total charge
\begin{equation}\label{eq:eight-upper}
 C\le35 f_0.
\end{equation}

On a $k$-gon, the payment is at least $5-k$: for triangles and
quadrilaterals this was proved, and for $k\ge5$ it follows from
nonnegativity. At a simple vertex in dimension 20 there are 190 two-faces, so
\[
 C\ge5f_2-190f_0.
\]
The nonnegative symmetric dual $h$-vector satisfies
$f_0=\sum h_i$ and $f_2=\sum\binom{20-i}{2}h_i$.
Pairing coefficients by $h_i=h_{20-i}$ gives the exact identity
\begin{equation}\label{eq:eight-gap}
 5f_2-225f_0
 =\frac52\sum_{i=0}^{20}(i-10)^2 h_i\ge500,
\end{equation}
because $h_0=h_{20}=1$. Hence $C\ge35f_0+500$, contradicting
\eqref{eq:eight-upper}. This excludes dimension 20 in $\mathcal K$.
Together with Theorem~\ref{thm:main}, it proves
Theorem~\ref{thm:eight19}.

The rank 20 step is written mathematics, not a new finite enumeration
or optimizer certificate. The source
\path{formal/CoxeterProofs/ArithmeticEightField19.lean} has compiled
and passed a separate printed-axiom audit of seven final theorems.
It checks the paired coefficients, endpoint surplus 500, quarter-scaled
cost update, elementary triangle/quad logic, the charge contradiction
and composition with the previous ceiling 20. Its printed axioms are
only \texttt{propext} and \texttt{Quot.sound}; the coefficient identity
is axiom-free. Actual geometry, local consumer classifications, complete
forest coverage and the identification of actual $h$-vector/charge data
remain explicit premises. This conditional kernel check is not an
end-to-end formalization of the geometric supplement.

\paragraph{Named internal supplement.}
The complete main consumer source is
\path{results/arithmetic_rank20_two_component_d_fork_half_refund_20261019.md};
the prior whole-forest theorem is
\path{results/arithmetic_rank20_complete_fourteen_family_frontier_20261019.md}.
For the local theorem Proposition~\ref{prop:eight-six}, the full
proofs are the following sources, all under \path{results/}:
\begin{longtable}{p{.13\linewidth}p{.78\linewidth}}
\toprule
Row&Complete source\\\midrule
I&\path{arithmetic_rank20_isolated_d_spectator_contact_obstruction_20261019.md}\\
F finite; T&\path{arithmetic_rank20_simple_d_finite_beta_and_far_fork_closure_20261019.md}\\
B finite&\path{arithmetic_rank20_far4_finite_beta_obstruction_20261019.md}\\
F/B large&\path{arithmetic_rank20_far_simple_and_far4_large_beta_closure_20261019.md}\\
J&\path{arithmetic_rank20_spectator_only_quad_obstruction_20261019.md}\\
J correction&\path{arithmetic_rank20_spectator_only_e8_coverage_addendum_20261019.md}\\
Bridge&\path{arithmetic_rank20_bridge_quad_obstruction_20261019.md}\\\bottomrule
\end{longtable}
The I composition explicitly imports its no-long, finite-long and
dotted-long unit proofs; their source pins and subsidiary arguments
are included in the supplement ledger. Its whole-spectator finite
continuation table is
\path{arithmetic_rank20_fork_spectator_continuation_research_20261019.md}.
The ledger also identifies every inherited prior charge theorem,
the separate written reviews and the conditional kernel evidence.

The source ledger distinguishes those external structural inputs,
inherited internal charge theorems, new local internal lemmas and the
conditional kernel records. To make this appendix self-contained at
the geometric level would require integrating all imported no-long,
finite-long and dotted-long norm arguments and the preceding $A/B/D$
consumer proofs, rather than replacing their source references by
review counts. No external refereeing, worldwide priority, sharpness
or new polytope example follows from the present theorem.
